\section{Preliminarii și context științific}

Studiul datelor colectate prin \emph{eye-tracking} se bazează pe ipoteza \emph{eye-mind}, conform căreia zonele în care privirea se oprește oferă informații despre procesele cognitive la momentul respectiv \cite{just1980}. Din punct de vedere tehnic, un traseu vizual (\emph{gazepath}) nu funcționează ca un flux continuu, ci mai degrabă este o alternanță între \emph{fixații} (acele perioade de timp în care ochiul rămâne relativ nemișcat pentru a observa detaliile dintr-o regiune) și \emph{sacade} (salturile rapide dintre două zone de interes, în timpul cărora vederea este suprimată). Într-o sarcină de căutare vizuală, precum \emph{Find Waldo}, înțelegerea și separarea acestor evenimente sunt necesare pentru a realiza o analiză corectă. Pentru a analiza matematic comportamentul participanților, trebuie efectuată trecerea de la o secvență continuă de coordonate la o secvență de stări discrete. Doar astfel putem măsura suprafața acoperită de privire, ordinea în care sunt inspectate zonele ecranului și interacțiunea directă cu personajul căutat.

Dincolo de extragerea acestor evenimente, un aspect important în studiul căutării vizuale este înțelegerea modului în care participanții își folosesc atenția pentru a rezolva sarcinile impuse. În acest sens, facem distincția între două tipuri de mecanisme: o căutare bazată pe reflexul vizual, numită \emph{bottom-up}, și căutarea orientată pe scopul sarcinii, numită \emph{top-down}. Modelele ce analizează caracteristicile vizuale ale unei scene (\emph{saliency-maps}), precum algoritmul clasic Itti-Koch, calculează elemente ce determină mecanismul \emph{bottom-up}, pornind strict de la imagine, precum intensitatea (contrastul), culoarea și orientarea contururilor (de exemplu, o muchie ascuțită sau o formă geometrică poate atrage privirea) \cite{itti1998}. Cu toate acestea, atunci când există un obiectiv clar pentru participanți, cum este găsirea personajului, intenția de a rezolva sarcina ajunge să controleze mișcările ochilor mult mai mult decât reflexele vizuale \cite{henderson2018meaning}. Pentru a putea măsura și evalua acest fenomen, trebuie să comparăm modul efectiv în care participanții au privit scenele cu hărțile de \emph{saliency} generate. Pentru ca această comparație să fie corectă, trebuie să folosim metrici de evaluare precum \emph{Normalized Scanpath Saliency} sau \emph{Information Gain}. Mai mult, rezultatele obținute trebuie raportate la un model de referință de bază, cum este \emph{center bias}, ce reprezintă tendința naturală a oamenilor de a privi direct spre centrul ecranului, indiferent de ce fel de imagine este afișată \cite{bylinskii2019metrics}. 

O altă direcție relevantă este aplicarea teoriei informației pentru a înțelege mai bine explorarea vizuală. Numărul de fixații nu este suficient pentru a descrie o strategie de căutare, deoarece două persoane pot avea același număr de fixații, dar trasee vizuale foarte diferite. De aceea, sunt utile conceptele de entropie spațială și entropie de tranziție. Cea spațială arată cât de răspândite sunt fixațiile pe imagine, iar cea de tranziție arată cât de previzibile sunt trecerile dintre zonele ecranului. Aceste metrici sunt potrivite pentru o sarcină de tip \emph{Find Waldo}, unde dificultatea nivelului poate schimba modul în care participantul explorează scena \cite{otero2016waldo}.

Nu în ultimul rând, o abordare tot mai folosită este trecerea de la simpla analiză statistică către modelarea prin \emph{Machine Learning}. Traseele vizuale sunt tratate fie ca serii de timp, fie ca date secvențiale. O direcție foarte utilă este prezicerea următoarei regiuni spre care se va uita un participant, folosind modele bazate pe lanțuri Markov, lucru care ne ajută să anticipăm pas cu pas mișcările următoare ale privirii. O altă abordare presupune folosirea rețelelor recurente (precum \emph{LSTM}). Aceste rețele păstrează informația despre ordinea fixațiilor, ceea ce le permite să surprindă diferența dintre o căutare care se apropie treptat de personaj și una în care participantul rămâne blocat în zone irelevante ale imaginii. În final, aceleași date ne permit să facem distincții clare între participanți, pe baza comportamentului vizual unic al fiecăruia.

Aceste concepte formează baza pentru restul lucrării: metricile extrase din traseele vizuale vor fi folosite atât pentru analiza comportamentului de căutare, cât și pentru antrenarea și evaluarea modelelor pe datele colectate.